\chapter{Detección de Binarios Troyanizados y Rootkits}

Los \textit{rootkits} a nivel de espacio de usuario (\textit{userland rootkits}) operaban reemplazando los binarios estándar de administración (como \texttt{ps}, \texttt{ls}, \texttt{netstat}, \texttt{find} y \texttt{login}) por versiones adulteradas que filtraban deliberadamente los procesos del atacante, ocultaban directorios maliciosos y abrían puertas traseras protegidas por contraseñas predefinidas.

\section{Búsqueda y análisis de programas troyanizados usando antivirus}

 
Muchos programas revisados muestran información y cadenas no
convencionales por esto se analizó toda la máquina con un
antivirus. Los archivos que contienen virus para windows se omiten del
reporte, pues es común que los usuarios los tengan en sus cuentas al llegar 
en correos o en programas bajados desde internet. En
linux los virus no son importantes y lo que se pretende es buscar
archivos con troyanos y gusanos hechos para linux y que pueden ser detectados 
por un antivirus.\\



 Ejemplos ilustrativos de los resultados de la revisión hecha por \verb+f-prot+ 
a 181,711 archivos y programas:
 
\vspace{1em}
\vspace{0.8em}\hrulefill\vspace{0.8em}
{\footnotesize
\begin{lstlisting}[style=terminal]
/root/usr/bin/last  Infection: Unix/RST.A
/root/usr/bin/strace  Infection: Unix/RST.A
/root/sbin/killall5  Infection: Unix/RST.A
/root/sbin/sulogin  Infection: Unix/RST.A
/root/bin/ps  Infection: Unix/RST.A
/root/dev/zqft4/local!  Infection: Unix/Osf.A
/home/staff/duende/paw++  Infection: Unix/RST.B
/home/staff/duende/prog/c/lc/une  Infection: Unix/RST.B
/home/staff/duende/prog/c/lc/linfo  Infection: Unix/RST.A
/home/staff/duende/prog/c/lc/mapea  Infection: Unix/RST.A
/home/staff/duende/prog/c/lc/particion  Infection: Unix/RST.A
/home/staff/duende/prog/c/lc/correlacion  Infection: Unix/RST.A
/home/staff/duende/prog/c/uno/cadena  Infection: Unix/RST.A
/home/staff/duende/prog/c/uno/arbol  Infection: Unix/RST.A
/home/staff/duende/prog/c/uno/arboles  Infection: Unix/RST.A
/home/staff/duende/prog/c/uno/repite  Infection: Unix/RST.A
/home/staff/duende/prog/c/uno/simotu  Infection: Unix/RST.A
/home/staff/duende/prog/GL/avion  Infection: Unix/RST.A
/home/staff/duende/prog/bin/une  Infection: Unix/RST.A
/home/staff/duende/prog/bin/normal  Infection: Unix/RST.A
/home/staff/duende/prog/gtk/ghello  Infection: Unix/RST.A
/home/staff/duende/prog/gtk/registro.old  Infection: Unix/RST.A
/home/staff/duende/kxterm  Infection: Unix/RST.A
/home/staff/duende/installer-redhat-70-i386  Infection: Unix/RST.A



Results of virus scanning:

Files: 190937
MBRs: 0
Boot sectors: 0
Objects scanned: 181711
Infected: 1111
Suspicious: 4
Disinfected: 0
Deleted: 0
Renamed: 0

Time: 55:33
\end{lstlisting}
}
\vspace{0.8em}\hrulefill\vspace{0.8em}
\vspace{1em}

De estos resultados, se encontraron cerca de 900 programas
troyanizados con gusanos que abren conexiones de red y dan acceso a
la máquina en diversos puertos. Fueron encontrados 115 programas infectados
en el directorio de duende principalmente en carpetas
que supuestamente contienen trabajo de investigación. Las fechas de
modificación de estos programas son de abril 24 pero los fuentes
correspondientes tienen fechas del 2001. Casi todos los programas esenciales en
los directorios de sistema están troyanizados. Hay que destacar que
las infecciones en linux no son automáticas y no se puede decir que
se infectaron solos por contacto. Los programas fueron infectados con cadenas
binarias de forma intencional y con claro acceso a la cuenta de
root. En un sistema limpio no debe aparecer ningún programa como
infectado como se verificó en los sistemas de control.



\section{Atributos especiales en archivos del sistema}


\vspace{1em}
\vspace{0.8em}\hrulefill\vspace{0.8em}
{\footnotesize
\begin{lstlisting}[style=terminal]
-u--ia----------- /usr/lib/popclean
-u--ia----------- /usr/local/lib/dsniff.services
-u--ia----------- /etc/cron.daily/dnsquery
-u--ia----------- /var/lib/.identd/libs
-u--ia----------- /var/lib/.identd/identd
-u--ia----------- /var/lib/.identd/libdn
-u--ia----------- /var/lib/.identd/libsnd
-u--ia----------- /dev/rd/c2d3b2
-u--ia----------- /dev/rd/c6d0b2
-u--ia----------- /dev/rd/c6d3b2
-u--ia----------- /dev/.libn
-u--ia----------- /dev/.libpt
-u--ia----------- /dev/.libtr
-u--ia----------- /dev/.libtp
-u--ia----------- /dev/.pst
-u--ia----------- /dev/.nst
-u--ia----------- /dev/.tree
-u--ia----------- /dev/.tp

\end{lstlisting}
}
\vspace{0.8em}\hrulefill\vspace{0.8em}
\vspace{1em}


El atributo \textbf{i} indica que los archivos anteriores no pueden ser
modificados o borrados por nadie. Estos atributos son conocidos por
muy pocas personas y un atacante no experimentado por lo general no
los sabe usar, pero son comunes cuando alguien instala un \textbf{rootkit}
empaquetado o compilado en el sistema que se esta atacando.

